\bigskip\Needspace{10\baselineskip}
\pdfbookmark[1]{Ceres Ferdinandea. 1802 November}{ceres180211}
\begin{center}
CERES FERDINANDEA.\par
\medskip\hrule\medskip
Monatliche Correspondenz. Band VI. S.\ 492—498. 1802 November.\par
\medskip\hrule
\end{center}\medskip
—\quad —\quad —

Dr.\ \textsc{Gauß} hat inzwischen weiter an der Bahn dieses Planeten gearbeitet. Unter Berücksichtigung der Störungen, wie er sie (Oktoberheft, S.\ 387) [S.\ 224 dieses Bandes] nach seinen Elementen VII berechnet hatte, zusammen mit den Breitengleichungen, bei denen seine Rechnung mit der von \textsc{Oriani} (Juniheft, S.\ 586. Juliheft, S.\ 68) [1802. Juli S.\ 67—70. —\quad —\quad —\quad \textsc{Oriani}, ein geschickter und in den verwickelten Störungsrechnungen so versierter Rechner, scheute keine Mühe, die gesamte Störungsrechnung für \textit{Ceres} unter einer anderen Annahme über die mittlere Entfernung zu wiederholen und damit zugleich seine zuerst erhaltenen Ergebnisse erneut zu prüfen. Hier ist das Ergebnis dieser ganzen Rechnung, die uns dieser große Astronom freundlicherweise mitgeteilt hat.
\[
\begin{aligned}
\text{Es sei }D&=\text{mittlere Länge der }\CeresSymbol-\text{mittlere Länge des }\JupiterSymbol;\\
A'&=\text{mittlere Anomalie }\JupiterSymbol;\\
A&=\text{mittlere Anomalie der }\CeresSymbol.\\
H'&=\text{mittlere Länge }\JupiterSymbol-\text{mittlere Länge }\AscendingNode\JupiterSymbol.\\
H&=\text{mittlere Länge }\CeresSymbol-\text{mittlere Länge }\AscendingNode\CeresSymbol.
\end{aligned}
\]
Ferner ist zu beachten, dass man die Störungsgleichungen für jede andere Exzentrizität $e$ der Bahn erhält, indem man die Glieder mit $A$ mit $\dfrac{e}{0.081406}$ multipliziert, die Glieder mit $2A$ mit $\left(\dfrac{e}{0.081406}\right)^2$ und die Glieder mit $3A$. mit $\left(\dfrac{e}{0.081406}\right)^3$.

\Needspace{24\baselineskip}
Es folgen die Gleichungen für die Störungen der \textit{Ceres} durch Jupiter und ihre Änderungen.
\SourceLeaf{233}%
Die folgende Tabelle gibt die Störungen der heliozentrischen Breite an.
\begin{center}
\begingroup\small
\begin{tabular}{l|r}
\shortstack{Nach den Bahnelementen VII\\von \textsc{Gauß}}&\shortstack{Wenn die mittlere jährliche\\Bewegung der Ceres um\\20 Min. erhöht wird.}\\\hline
$-13''19\sin(H-D)$&$-12''97$\\
$+1.62\sin H'$&$+1.59$\\[3pt]
$+16.21\sin(2D-H)$&$+16.03$\\
$-1.99\sin(D-H')$&$-1.96$\\[3pt]
$+32.38\sin(3D-H)$&$+32.78$\\
$-3.97\sin(2D-H')$&$-4.01$\\[3pt]
$-5.50\sin(4D-H)$&$-5.31$\\
$+0.67\sin(3D-H')$&$+0.65$\\[3pt]
$+6.17\sin(D+H)$&$+6.06$\\
$-0.76\sin(2D+H')$&$-0.74$\\[3pt]
$+1.52\sin(2D+H)$&$+1.49$\\
$-0.19\sin(5D+H')$&$-0.18$\\[3pt]
$+16.15\sin(4D-A'-H-2^0\ 33')$&$+13.83$\\
$-13.86\sin(5D-A-H-2^0\ 33')$&$-11.92$
\end{tabular}\endgroup
\end{center}
fast vollständig übereinstimmt], bestimmte er sogleich die folgenden Elemente.
\begin{center}
\begingroup\small\setlength{\tabcolsep}{3pt}
\begin{tabular}{l r}
Epoche 1801 für Seeberg $\ .\ .$&$77^0\ 19'\ 38''4$\\
Tägliche tropische Bewegung $\ .$&$770''764$\\
Umlaufszeit $\ .\ .\ .\ .\ .\ .\ .$&1681 Tage 11 Stunden\\
Exzentrizität $\ .\ .\ .\ .\ .\ .\ .$&$0.0788132$\\[3pt]
Logarithmus der großen Halbachse $\ .$&$0.4421085$\\
Aphel 1801 $\ .\ .\ .\ .\ .\ .\ .\ .\ .$&$326^0\ 33'\ 10''$\\
Knoten 1801 $\ .\ .\ .\ .\ .\ .\ .\ .\ .\ .$&$80\ 54\ 52$\\
Neigung — $\ .\ .\ .\ .\ .\ .\ .\ .\ .\ .$&$10\ 37\ 48$
\end{tabular}\endgroup
\end{center}

Mit diesen Elementen berechnete Dr.\ \textsc{Gauß} nun die Störungen der $\CeresSymbol$ durch $\JupiterSymbol$ erneut. Die Ergebnisse folgen.
\begin{center}\textit{Säkulare Gleichungen.}\end{center}
\begin{center}
\begingroup\small\setlength{\tabcolsep}{3pt}
\begin{tabular}{l r}
Jährliches Vorrücken des Aphels gegenüber den Äquinoktien&$+121''3$\\
Jährliches Vorrücken der Knoten gegenüber den Äquinoktien $\ .\ .$&$+0.8$\\
Jährliche Abnahme der Exzentrizität $\ .\ .\ .\ .\ .\ .\ .\ .\ .$&$0.000005829$\\
Jährliche Abnahme der Bahnneigung $\ .\ .\ .\ .\ .\ .\ .$&$+0''38$
\end{tabular}\endgroup
\end{center}
\Needspace{19\baselineskip}
\begin{center}\textit{Periodische Gleichungen.}\par\medskip
\textbf{In der Länge.}\end{center}
\par\noindent\begin{minipage}[t]{.48\linewidth}\small
\[
\begin{aligned}
&-231''00\sin(\CeresSymbol-\JupiterSymbol)\\
&+492.37\sin2(\CeresSymbol-\JupiterSymbol)\\
&+43.85\sin3(\CeresSymbol-\JupiterSymbol)\\
&+10.00\sin4(\CeresSymbol-\JupiterSymbol)\\
&+3.03\sin5(\CeresSymbol-\JupiterSymbol)\\
&+1.04\sin6(\CeresSymbol-\JupiterSymbol)\\
&+0.40\sin7(\CeresSymbol-\JupiterSymbol)\\
&+0.16\sin8(\CeresSymbol-\JupiterSymbol)\\
&-58.90\sin(\JupiterSymbol+16^0\ 56'\ 37'')\\
&-598.69\sin(\CeresSymbol-2\JupiterSymbol-26^0\ 6'\ 58'')\\
&-437.75\sin(2\CeresSymbol-3\JupiterSymbol-10^0\ 54'\ 25'')
\end{aligned}
\]
\end{minipage}\hfill\begin{minipage}[t]{.515\linewidth}\small
\[
\begin{aligned}
&+54''90\sin(3\CeresSymbol-4\JupiterSymbol-10^0\ 32'\ 23'')\\
&+10.28\sin(4\CeresSymbol-5\JupiterSymbol-10\quad21\quad35'')\\
&+3.19\sin(5\CeresSymbol-6\JupiterSymbol-10\quad15\quad2'')\\
&+1.20\sin(6\CeresSymbol-7\JupiterSymbol-10\quad10\quad20'')\\
&+0.50\sin(7\CeresSymbol-8\JupiterSymbol-10\quad6\quad24'')\\
&+22.72\sin(2\CeresSymbol-\JupiterSymbol+36\quad0\quad12'')\\
&-51.80\sin(3\CeresSymbol-2\JupiterSymbol+32\quad47\quad57'')\\
&-5.74\sin(4\CeresSymbol-3\JupiterSymbol+30\quad41\quad32'')\\
&-1.65\sin(5\CeresSymbol-4\JupiterSymbol+28\quad53\quad58'')\\
&-0.69\sin(6\CeresSymbol-5\JupiterSymbol+28\quad13\quad52'')\\
&-0.26\sin(7\CeresSymbol-6\JupiterSymbol+26\quad41\quad16'')
\end{aligned}
\]
\end{minipage}\par
\SourceLeaf{234}%
\Needspace{19\baselineskip}
\begin{center}\textbf{Im Radiusvektor.}\end{center}
\par\noindent\begin{minipage}[t]{.48\linewidth}\small
\[
\begin{aligned}
&-0.0000943\\
&+0.0010261\cos(\CeresSymbol-\JupiterSymbol)\\
&-0.0037679\cos2(\CeresSymbol-\JupiterSymbol)\\
&-0.0004176\cos3(\CeresSymbol-\JupiterSymbol)\\
&-0.0001069\cos4(\CeresSymbol-\JupiterSymbol)\\
&-0.0000347\cos5(\CeresSymbol-\JupiterSymbol)\\
&-0.0000127\cos6(\CeresSymbol-\JupiterSymbol)\\
&-0.0000050\cos7(\CeresSymbol-\JupiterSymbol)\\
&-0.0000021\cos8(\CeresSymbol-\JupiterSymbol)\\
&-0.0000614\cos(\CeresSymbol+24^0\ 13'\ 37'')\\
&+0.0002398\cos(\JupiterSymbol+22\quad58\quad33'')\\
&+0.0008325\cos(\CeresSymbol-2\JupiterSymbol-23^0\ 54'\ 51'')
\end{aligned}
\]
\end{minipage}\hfill\begin{minipage}[t]{.515\linewidth}\small
\[
\begin{aligned}
&+0.0025560\cos(2\CeresSymbol-3\JupiterSymbol-10^0\ 39'\ 20'')\\
&-0.0004698\cos(3\CeresSymbol-4\JupiterSymbol-10\quad51\quad24'')\\
&-0.0001048\cos(4\CeresSymbol-5\JupiterSymbol-10\quad46\quad51'')\\
&-0.0000358\cos(5\CeresSymbol-6\JupiterSymbol-10\quad40\quad14'')\\
&-0.0000142\cos(6\CeresSymbol-7\JupiterSymbol-10\quad34\quad4'')\\
&-0.0000061\cos(7\CeresSymbol-8\JupiterSymbol-10\quad28\quad56'')\\
&-0.0001313\cos(2\CeresSymbol-\JupiterSymbol+37\quad11\quad18'')\\
&+0.0002928\cos(3\CeresSymbol-2\JupiterSymbol+32\quad17\quad53'')\\
&+0.0000397\cos(4\CeresSymbol-3\JupiterSymbol+29\quad3\quad25'')\\
&+0.0000135\cos(5\CeresSymbol-4\JupiterSymbol+26\quad58\quad26'')\\
&+0.0000057\cos(6\CeresSymbol-5\JupiterSymbol+25\quad46\quad37'')\\
&+0.0000026\cos(7\CeresSymbol-6\JupiterSymbol+25\quad2\quad49'')
\end{aligned}
\]
\end{minipage}\par
\begin{center}\textbf{In der Breite.}\end{center}
\par\noindent\begin{minipage}[t]{.48\linewidth}\small
\[
\begin{aligned}
&-11''60\sin(\JupiterSymbol-78^0\ 31'\ 50'')\\
&+14.29\sin(\CeresSymbol-2\JupiterSymbol+78^0\ 31'\ 50'')\\
&+28.73\sin(2\CeresSymbol-3\JupiterSymbol+78\quad31\quad50'')\\
&-4.81\sin(3\CeresSymbol-4\JupiterSymbol+78\quad31\quad50'')\\
&-0.96\sin(4\CeresSymbol-5\JupiterSymbol+78\quad31\quad50'')
\end{aligned}
\]
\end{minipage}\hfill\begin{minipage}[t]{.515\linewidth}\small
\[
\begin{aligned}
&-0''29\sin(5\CeresSymbol-6\JupiterSymbol+78^0\ 31'\ 50'')\\
&+5.43\sin(2\CeresSymbol-\JupiterSymbol-78\quad31\quad50'')\\
&+1.34\sin(3\CeresSymbol-2\JupiterSymbol-78\quad31\quad50'')\\
&+0.44\sin(4\CeresSymbol-3\JupiterSymbol-78\quad31\quad50'')
\end{aligned}
\]
\end{minipage}\par
Schließlich berechnete Dr.\ \textsc{Gauß} mit diesen Formeln die Zahlenwerte der Störungen für die zugrunde gelegten Beobachtungen erneut. So fand er die folgenden Bahnelemente, die wir künftig mit (VIII) bezeichnen werden.
\begin{center}
\begin{tabular}{r|r|r|r}
\multicolumn{2}{c|}{Epoche für Seeberg}&Aphel&Knoten\\\hline
1801&$77^0\ 19'\ 34''9$&$326^0\ 33'\ 37''$&$80^0\ 54'\ 59''$\\
1802&$155\ 28\ 35.1$&$326\ 35\ 39$&$80\ 55\ 0$\\
1803&$233\ 37\ 35.3$&$326\ 37\ 40$&$80\ 55\ 1$
\end{tabular}
\end{center}
\begin{center}
\begingroup\small\setlength{\tabcolsep}{3pt}
\begin{tabular}{l r}
Tägliche tropische Bewegung $\ .\ .\ .\ .\ .\ .\ .\ .\ .\ .\ .$&$770''7951$\\
Tropische Umlaufszeit $\ .\ .\ .\ .\ .\ .\ .\ .\ .\ .\ .\ .\ .$&1681 Tage 9 Stunden\\
Exzentrizität (1801) $\ .\ .\ .\ .\ .\ .\ .\ .\ .\ .\ .\ .\ .\ .$&$0.0788352$\\
Neigung (1801) $\ .\ .\ .\ .\ .\ .\ .\ .\ .\ .\ .\ .\ .\ .\ .\ .$&$10^0\ 37'\ 56''0$\\
Logarithmus der Halbachse $\ .\ .\ .\ .\ .\ .\ .\ .\ .\ .$&$0.4420971$
\end{tabular}\endgroup
\end{center}
Diese neuen Elemente unterscheiden sich so wenig von den vorhergehenden, dass man eine erneute Berechnung der Störungen nach ihnen als überflüssig ansehen kann, da die Ergebnisse von den oben angegebenen kaum merklich abweichen können. Aus Neugier prüfte Dr.\ \textsc{Gauß}, wie stark die neuen Elemente VIII um die Zeit des Zusammentreffens im kommenden Jahr von den Elementen VII abweichen werden. Ohne die Aberration zu berücksichtigen, findet er für den 28. Juni 1803 um 12 Uhr mittlerer Seeberger Zeit die folgenden Werte.
\begin{center}
\begin{tabular}{c|r|r}
\shortstack{Nach den\\Elementen}&Länge der $\CeresSymbol$&\shortstack{Breite der $\CeresSymbol$\\südl.}\\\hline
VII&$280^0\ 7'\ 43''$&$5^0\ 1'\ 21''$\\
VIII&$280\ 17\ 58$&$5\ 4\ 53$\\\hline
Unterschied&$10\ 15$&$3\ 32$
\end{tabular}
\end{center}
\SourceLeaf{235}%

Dr.\ \textsc{Gauß} schließt seine Rechnung mit dieser Bemerkung. „Da es schon nach den obigen Störungsgleichungen so viel Mühe macht, eine Position des Planeten zu berechnen, scheint es wohl noch kaum ratsam, ihre Zahl bereits jetzt durch die Gleichungen für Saturn und diejenigen für Jupiter zu erhöhen, die von den Quadraten und Produkten der Exzentrizitäten abhängen. Ihre Aufnahme böte im Übrigen keine weitere Schwierigkeit als das \textit{Taedium} des mechanischen Rechnens, zumal wir den Elementen jetzt nach fast einem halben Jahr mithilfe neuer Beobachtungen eine viel größere Genauigkeit werden geben können. Da es so mühsam ist, den jeweiligen Betrag der Gleichungen zu berechnen, wird es meines Erachtens wichtig sein, über eine Kürzung und Vereinfachung ihrer Tafeln nachzudenken. Ich werde Ihnen meine Gedanken dazu später zusammen mit einem \textit{Specimen} vorlegen.“
\begin{flushright}[\textsc{von Zach.}]\end{flushright}

